\section{Related Work}
\label{sec:related_work}
The proliferation of deep learning in industrial fault diagnosis has addressed many challenges associated with manual feature extraction. However, the performance of these models remains highly dependent on the availability of balanced and high-quality labeled datasets. This section reviews the evolution of Generative Adversarial Networks (GANs) in mechanical engineering, the advancements in domain translation using CycleGAN, and existing strategies for imbalanced learning in fault diagnosis.

\subsection{GANs in Fault Diagnosis}
Since their introduction by Goodfellow \cite{1}, Generative Adversarial Networks (GANs) have revolutionized the field of synthetic data generation by pitting a generator against a discriminator in a minimax game. In the context of mechanical fault diagnosis, early GAN variants were often plagued by training instabilities and mode collapse, where the generator produces a limited variety of samples. To mitigate these issues, Arjovsky et al. \cite{2} proposed the Wasserstein GAN (WGAN), which utilizes the Earth Mover's distance to provide a more stable loss function. This was further refined by Gulrajani et al. \cite{3} through the introduction of a Gradient Penalty (WGAN-GP), ensuring the Lipschitz constraint without the drawbacks of weight clipping.

In mechanical engineering, these architectures have been adapted to capture the complex temporal dependencies of vibration signals. For instance, the Conditional WGAN-GP (CWGAN-GP) \cite{7} allows for the generation of class-specific samples, which is crucial for augmenting minority fault classes. Extensive research has shown that integrating physical constraints into GANs can improve the diagnostic accuracy of downstream classifiers by providing more realistic training data. Despite these advancements, standard GAN architectures often struggle with high-fidelity signal reconstruction when the mapping between the latent space and the physical signal domain is not well-constrained. Most existing GAN-based methods focus on generating data from random noise, which may lack the structural consistency required for precise diagnosis in varying operating conditions. This limitation highlights the need for more structured generative approaches that can leverage existing signal domains.

\subsection{Domain Translation and CycleGAN}
Domain translation has emerged as a powerful paradigm for transforming data from one distribution to another while preserving essential features. The Cycle-Consistent Adversarial Network (CycleGAN), proposed by Zhu et al. \cite{4}, introduced a cycle-consistency loss that enables unpaired image-to-image translation. This concept has recently been extended to the 1D signal domain for fault diagnosis applications. For example, 1D-CycleGAN \cite{37} has been utilized to translate vibration signals across different motor speeds, effectively addressing the data scarcity problem under specific working conditions.

Compared to traditional data generation techniques like signal modeling or simple interpolation, CycleGAN-based translation \cite{36} offers a data-driven approach that captures the underlying physical transformation between healthy and faulty states. Recent studies \cite{36} have demonstrated that translating signals from a source domain (e.g., simulation or laboratory data) to a target domain (e.g., real-world industrial data) can significantly enhance the robustness of diagnostic classifiers. In many industrial scenarios, vibration signals from the same component under different loads or speeds share underlying fault characteristics, which CycleGAN can effectively exploit. However, many current domain translation methods in fault diagnosis do not explicitly account for the auxiliary classification information during the generation process. This often leads to a "semantic shift" where the translated fault features do not perfectly align with the intended fault category, necessitating a more integrated approach that combines domain translation with class-conditional guidance.

\subsection{Imbalanced Learning Strategies}
The class imbalance problem is a ubiquitous challenge in industrial fault diagnosis, as data from faulty states is naturally rarer than data from healthy states. Classical oversampling techniques such as Synthetic Minority Over-sampling Technique (SMOTE) \cite{6} and ADASYN \cite{35} have been widely used to balance datasets by interpolating between existing minority samples. While effective for low-dimensional data, these methods often fail to capture the complex, non-linear manifolds of high-frequency vibration signals, leading to the generation of "noisy" samples that overlap with majority classes. This overlap can significantly degrade the classification boundary of neural networks, leading to high false-negative rates for critical mechanical failures.

To address these shortcomings, hybrid deep learning approaches have been developed. The Auxiliary Classifier GAN (ACGAN) \cite{5} incorporates a classification branch to ensure the generated samples are class-discriminative. Building on this, researchers have proposed variants like ACWGAN \cite{35} and Improved Auxiliary Conditional WGAN-GP (IACWGAN-GP) \cite{26} to combine the stability of WGAN-GP with the class-steering capabilities of ACGAN. These hybrid models aim to produce samples that are both physically realistic and highly separable in the feature space. Despite these efforts, a gap remains in effectively combining the cross-domain mapping strengths of CycleGAN with the rigorous class-conditioning of AC-based frameworks. Our proposed method addresses this by integrating auxiliary classification loss into a cyclic translation framework to ensure both domain-level realism and category-level precision for imbalanced fault diagnosis. This dual-objective optimization strategy allows for the synthesis of highly accurate fault samples even in the absence of paired training data.
